\documentclass[11pt,a4paper]{article}
\usepackage[margin=1in]{geometry}
\usepackage{amsmath,amssymb,amsfonts,graphicx,booktabs,hyperref,url}
\usepackage{algorithm}
\usepackage{algorithmic}
\usepackage{subcaption}

\title{\textbf{RaZOOM: Sample-Efficient Embodied Control with Frozen Semantic Priors and Differentiable Physics}\\[0.4em]
\large Technical Report v1.0}
\author{Illia Ovcharenko \\ RaZOOM Research \\ \texttt{contact@razoom.ai}}
\date{May 2026}

\begin{document}
\maketitle

\begin{abstract}
Embodied reinforcement learning often depends on large human demonstration datasets or billion-parameter foundation policies, limiting deployment where teleoperation is unsafe or uneconomical. We present \textbf{RaZOOM}, combining differentiable physics (MuJoCo MJX) with a \textbf{multi-stream policy} that separates reflexive control, affect-modulated valuation, and topology-routed deliberation. A \textbf{frozen semantic prior} (${\sim}$2.6M parameters) supplies structured inductive bias while ${\sim}$157K policy parameters are updated during learning, enabling single-GPU training and CPU inference under 10\,ms. On NVIDIA H200 rollouts (${\ge}$345K environment steps in monolithic jobs) and controlled 100K-step A/B studies, agents reach task rewards ${\sim}$0.70 with stable homeostatic metrics. Optional sensory grounding trades throughput (${\sim}$159 vs.\ ${\sim}$108 env-steps/s) for comparable task success. We describe a sanitized evaluation protocol; proprietary prior databases are not released.
\end{abstract}

\section{Introduction}
\label{sec:intro}

Training physical agents at scale remains bottlenecked by \emph{data}, not compute. Industrial manipulation, hazardous material handling, and remote logistics often forbid the collection of large teleoperation corpora. Recent vision-language-action (VLA) and foundation policies achieve strong performance but require web-scale pretraining, cloud inference, and demonstration-heavy fine-tuning~\cite{brohan2023rt2,ye2024rlfp}. Model-based reinforcement learning (MBRL) with learned world models reduces sample complexity in simulation~\cite{hafner2023dreamerv3}, yet struggles on novel tasks without structural priors.

We argue that \textbf{structured inductive bias}---not parameter count---is the lever for sample-efficient embodied control. \textbf{RaZOOM} is a lightweight stack that couples (i) \textbf{differentiable physics simulation} (MJX) for gradient-friendly rollouts, (ii) a \textbf{multi-stream policy} with parallel reflex, affect, and planning pathways, and (iii) a \textbf{frozen semantic prior} encoding domain structure without online supervision. The trainable policy is compact (${\sim}$157K parameters); priors remain fixed (${\sim}$2.6M parameters), decoupling representation learning from control learning.

\paragraph{Contributions.}
\begin{itemize}
  \item A multi-stream embodied architecture with emotion-conditioned fast control and graph-topology-conditioned slow planning, integrated with canonical semantic embeddings.
  \item A JAX-native training pipeline (REINFORCE/PPO, JEPA-style regularization) over MJX with MLflow instrumentation and cloud orchestration (SkyPilot, H200).
  \item Empirical evidence from Nebius H200 campaigns: 100K-step A/B comparisons and long-horizon monolithic runs, including throughput--accuracy trade-offs for grounded sensing.
  \item A reproducibility package scope: sanitized APIs and metrics protocol; weights and prior corpora withheld for competitive reasons.
\end{itemize}

\paragraph{Results preview.}
In aggregated 100K-step studies, baseline agents achieve mean task reward ${\sim}$0.70 at ${\sim}$159 environment steps/s; grounded sensory variants match task reward at ${\sim}$108 steps/s (Table~\ref{tab:main100k}). Homeostatic metrics (conflict rate, nutrition gap) remain stable in JAX rollouts logged to MLflow.

\section{Related Work}
\label{sec:related}

\paragraph{Sample-efficient embodied RL.}
RL with foundation priors (RLFP) demonstrates that frozen value and success models can accelerate real-robot learning~\cite{ye2024rlfp}. RaZOOM differs by using \emph{structured symbolic priors} rather than general VLMs, targeting edge deployability.

\paragraph{World models and differentiable simulation.}
DreamerV3 learns latent dynamics for long-horizon control~\cite{hafner2023dreamerv3}. Differentiable engines (Brax/MJX) enable gradients through physics~\cite{freeman2021mjx}. We use MJX as the environment backend while learning a factored policy rather than a monolithic world model.

\paragraph{Hierarchical and affective control.}
Hierarchical RL decomposes tasks into sub-policies; homeostatic and intrinsic motivation literature motivates separate valuation streams~\cite{makoviychuk2021isaac}. Our valence stream implements affect-modulated gains on fast pathways, stabilizing exploration without hand-crafted reward shaping.

\paragraph{Joint-embedding predictive architectures.}
JEPA encourages representation structure via predictive losses~\cite{lecun2022jepa}. We apply a lightweight hierarchical regularization term during training ($\alpha{=}0.25$, $\lambda{=}0.1$) without releasing proprietary embedding corpora.

\section{Problem Formulation}
\label{sec:problem}

We model embodied control as a partially observable Markov decision process (POMDP) $(\mathcal{S}, \mathcal{A}, \mathcal{O}, T, R, \gamma)$. Observations $o_t \in \mathcal{O}$ concatenate proprioceptive state, optional exteroceptive features, and a low-dimensional \emph{configuration index} over a fixed topology graph. Actions factor into four branches: locomotion ($|\mathcal{A}_{\mathrm{loc}}|{=}16$), respiration ($2$), head ($2$), and macro ($4$), totaling 24 discrete decisions per step.

\paragraph{Homeostatic objective.}
Beyond extrinsic task reward $r^{\mathrm{task}}_t$, we log intrinsic homeostatic signals: food concentration, nutrition gap, conflict rate, and population-level coordination statistics. The training objective combines policy-gradient loss on $r^{\mathrm{task}}$ with auxiliary regularizers (JEPA and hierarchical structure terms). We do not assume access to human trajectories.

\paragraph{Parameter budget.}
Let $\theta$ denote trainable policy parameters ($|\theta|{\approx}1.57{\times}10^5$) and $\phi$ frozen semantic priors ($|\phi|{\approx}2.6{\times}10^6$). Only $\theta$ is updated during RL; $\phi$ is loaded from an offline corpus (not publicly released).

\section{Method}
\label{sec:method}

\subsection{System overview}

Figure~\ref{fig:architecture} illustrates RaZOOM. Per timestep, the agent encodes observations, retrieves a frozen semantic embedding, and routes decisions through three streams before actuating the differentiable simulator.

\begin{figure}[t]
  \centering
  \fbox{\parbox{0.92\linewidth}{\centering\vspace{2.8em}
  \textbf{[Figure 1: Multi-stream architecture]}\\
  \small Observations $\rightarrow$ Frozen semantic encoder $\rightarrow$ \\
  Fast (reflex) $\parallel$ Valence (affect) $\parallel$ Slow (topology route) $\rightarrow$ MJX\\
  \vspace{2.8em}}}
  \caption{RaZOOM control stack. Public diagrams omit proprietary graph connectivity and prior token identities.}
  \label{fig:architecture}
\end{figure}

\subsection{Frozen semantic prior}

A corpus of $N_{\mathrm{sym}}{=}110$ canonical symbols is embedded in $\mathbb{R}^{d_{\mathrm{emb}}}$ ($d_{\mathrm{emb}}{=}768$) offline. A projection head maps embeddings to policy features $z_t \in \mathbb{R}^{d_z}$. During RL, $\phi$ is frozen; gradients do not modify the corpus. This induces a \textbf{relational inductive bias}: the agent starts with structured semantics rather than learning language from scratch in simulation.

\subsection{Multi-stream policy}

Let $h_t$ denote concatenated proprioceptive/visual features. The policy factorizes as
\begin{equation}
\pi_\theta(a_t \mid h_t, z_t, e_t, \rho_t) \propto
\pi_{\mathrm{fast}}(a^{\mathrm{loc}}_t, a^{\mathrm{head}}_t \mid h_t, z_t, e_t)\,
\pi_{\mathrm{val}}(a^{\mathrm{breath}}_t, e_{t+1} \mid h_t, z_t, e_t)\,
\pi_{\mathrm{slow}}(a^{\mathrm{macro}}_t \mid h_t, z_t, \rho_t),
\end{equation}
where $e_t \in \mathbb{R}^{d_e}$ ($d_e{=}7$) is an affect state and $\rho_t \in \mathbb{R}^{d_\rho}$ is a route vector over a fixed topology graph ($d_\rho{=}22$).

\paragraph{Fast stream (reflex).}
Emotion-conditioned affine maps compute logits for locomotion and head actions:
\begin{equation}
\tilde{h}_t = h_t \odot \bigl(1 + \gamma \cdot \mathrm{mean}(e_t)\bigr), \quad
\ell^{\mathrm{loc}}_t = W_{\mathrm{loc}} \tilde{h}_t + b_{\mathrm{loc}},
\end{equation}
with learnable gain $\gamma$.

\paragraph{Valence stream.}
The valence stream updates $e_{t+1}$ and breath logits, modulating fast-stream gains. This implements \textbf{affect-modulated exploration}: high arousal increases sensitivity of reflex policies without manual reward engineering.

\paragraph{Slow stream (topology routing).}
A graph-topology-conditioned layer maps $(h_t, z_t, \rho_t)$ to macro-action logits. Route vectors index edges on a fixed planning graph with horizon $H{=}22$. We treat $\rho_t$ as a discrete--continuous hybrid index into a 64-state Markov conditioner (public name: \emph{discrete state enum}).

\subsection{Differentiable physics backend}

Rollouts execute in MuJoCo MJX with parallel agents (up to 16). Gradients may flow through smooth dynamics where supported; non-differentiable contact events use surrogate gradients consistent with MBRL practice. This backend supplies high-throughput data generation on H200 GPUs.

\subsection{Training and consolidation}

\paragraph{Online RL.}
We implement REINFORCE and PPO with generalized advantage estimation. Checkpoints serialize policy weights as portable JSON for auditability.

\paragraph{Regularization.}
Total loss:
\begin{equation}
\mathcal{L} = \mathcal{L}_{\mathrm{RL}} + \alpha \mathcal{L}_{\mathrm{JEPA}} + \lambda \mathcal{L}_{\mathrm{struct}},
\end{equation}
with $(\alpha,\lambda){=}(0.25, 0.1)$ from our training configuration.

\paragraph{Sleep replay consolidation.}
High-resonance trajectory blocks trigger shallow offline updates (policy nudge) without unfreezing $\phi$. This mimics offline experience replay while keeping compute bounded.

\subsection{Implementation notes}

RaZOOM is implemented in Python with JAX (no PyTorch). Orchestration uses SkyPilot on Kubernetes (Nebius H200). Metrics (steps per second, homeostasis, conflict rate) log to MLflow. We withhold source-level module names, graph construction rules, and prior databases.

\section{Experiments}
\label{sec:experiments}

\subsection{Infrastructure}

Training runs on NVIDIA H200 (MIG slices, CUDA 13) with containerized dependencies (\texttt{.venv-nebius-mjx-cuda13}). Object storage (endpoint \texttt{s3.nebius.ru}, bucket \texttt{nebius-pixlie-models}) persists checkpoints. Monolithic jobs target up to $1.5{\times}10^6$ environment steps; A/B studies use $10^4$--$10^5$ steps for controlled comparisons.

\subsection{Protocols}

\paragraph{A/B non-Dreamer compare.}
We compare \texttt{baseline} against \texttt{rich128\_grounded} (grounded sensory append) and sensory variants. Metrics: task reward, reward mean, intrinsic reward, approximate steps/s, checkpoint count.

\paragraph{JAX monolithic.}
\texttt{jax\_razoom\_runtime.py} logs SPS, homeostasis variance, conflict rate, food concentration, nutrition gap, and discrete state index.

\paragraph{Baselines.}
We report DreamerV3 configurations present in the repository as exploratory baselines; primary claims use RaZOOM multi-stream runs.

\subsection{Evaluation metrics}

Task success (task reward), aggregate reward mean, throughput (environment steps/s), and homeostatic scalars. We aggregate across seeds where available and report means in Table~\ref{tab:main100k}.

\input{tables_main.tex}

\begin{figure}[t]
  \centering
  \fbox{\parbox{0.92\linewidth}{\centering\vspace{1.5em}
  \textbf{[Figure 2: Throughput vs.\ task reward]}\\
  \small See \texttt{research\_paper\_resources/figures/fig2\_throughput\_vs\_task\_reward.svg}
  \vspace{1.5em}}}
  \caption{Throughput vs.\ task reward across Nebius compare runs (38 logged A/B jobs; marker size scales with training steps).}
  \label{fig:scatter}
\end{figure}

\section{Results and Discussion}
\label{sec:results}

\paragraph{Throughput--accuracy trade-off.}
Grounded sensory variants reduce throughput (${\sim}$108--136 steps/s) relative to baseline (${\sim}$155--163 steps/s) while maintaining task reward ${\sim}$0.70 (Figure~\ref{fig:scatter}, Table~\ref{tab:main100k}). This supports using grounding when sample efficiency matters more than raw rollout speed.

\paragraph{Long-horizon trends.}
Baseline task reward remains near 0.70 for horizons $10^4$--$10^5$ steps; at $2{\times}10^5$ steps some runs show lower task reward (${\sim}$0.56--0.70), motivating future curriculum and replay tuning (Appendix).

\paragraph{Deployment.}
Policy inference targets CPU execution under 10\,ms per decision for edge deployment; exact latency depends on observation dimensionality (benchmarked separately).

\paragraph{Comparison to foundation approaches.}
Unlike billion-parameter VLA models, RaZOOM prioritizes \textbf{structure over scale}: frozen priors plus compact $\theta$ enable training in hours on a single GPU rather than weeks on clusters.

\section{Limitations}
\label{sec:limitations}

\begin{itemize}
  \item \textbf{Simulation-only evidence.} Primary results are from MJX; sim-to-real transfer via Unity/ROS2 is in progress.
  \item \textbf{Withheld priors.} Reproducibility is limited to protocol and sanitized APIs; $\phi$ is proprietary.
  \item \textbf{Incomplete ablations.} Some ablation cells require additional short reruns (frozen prior removal, slow-stream removal).
  \item \textbf{Metric definitions.} Intrinsic reward can be negative by construction; readers should interpret signs relative to homeostatic design.
\end{itemize}

\section{Conclusion}
\label{sec:conclusion}

RaZOOM demonstrates that embodied agents can train without human demonstrations when structured semantic priors and differentiable physics are combined with a multi-stream policy. Future work targets real-robot pilots, public inference APIs, and open benchmarks with sanitized priors.

\bibliographystyle{plain}
\bibliography{references}

\appendix
\input{appendix.tex}

\end{document}
